\scanpage{61}
\mspage{60}

Pratiquement, quand on ne travaille pas avec des ``lieux'' proprement dits (et je suis le premier, apparemment, à y avoir songé, aux lieux\ldots) mais avec des points, et dans l'optique des espaces ``infiniment divisibles'' (et avec \unsure{bouzy|Bouzy|Bourb.}, je suis le premier à songer à travailler \unsure{éo!|eô!|eo!} avec des espaces qui ne le soient pas), il semblerait que \deleted{\unsure{ça soit bon}} cette structure \unsure{satisfasse|satisfera} à tous les besoins (avec, bien sûr, des axiomes \uline{en plus}, suivant les besoins). Cette structure a l'avantage d'être \uline{commune} aux contextes ``\unsure{spacieux|spécieux}'' et ``polyédral'', et se recommande de plus par la simplicité étonnante des \uline{données} (seulement $\leq$ et $\ll$ dans $M$ \unsure{lui-même} dans le cas des ateliers $F$), et des axiomes, réduits ici au nb de quatre (à peine plus compliqué qu'une topologie!). Quand on travaille avec des \deleted{\unsure{Quand}} ateliers à figures finies\inserted{ (cas le plus important)}, ceux-ci sont définis déjà par le préatelier correspondant.

% margin note written sideways; slanted upward, crossing the left edge of the paragraph
\margin{Ça mériterait peut-être un nom plus lapidaire que ``Préatelier quasi-ensembliste'', p.ex. ``\unsure{Stratop}'', ou (avec moins d'emphase) ``Stratologie'' (à moins d'utiliser ``Préatelier'' sans plus). L'usage dira!}

Si on veut travailler avec des ateliers qu. ens. (\unsure{spacieux|spécieux} resp. polyédraux) plus généraux, il faut introduire sur $M$ la notion de \deleted{\unsure{famille}} partie (fermée) ``admissible'' par des conditions convenables de ``locale finitude''.
